Le laiton est un alliage de cuivre et de zinc. Quand cet alliage est mis en
présence d'une solution d'acide chlorhydrique, seul le zinc réagit~; il se forme
du dihydrogène et des ions \ch{Zn^{2+}} en solution.

On attaque $\qty{10}{\g}$ de cet alliage par $\qty{100}{\mL}$ de solution
d'acide chlorhydrique dont la concentration est $c=\qty{1.00}{\mol\per\L}$. Il
se forme alors $V=\qty{0.900}{\L}$ de dihydrogène, mesuré dans les conditions
normales de température et de pression ($V_{m}=\qty{22.4}{\L\per\mol}$).
\begin{enumerate}
    \item Écrire l'équation de la réaction qui a lieu.
    \item Calculer la quantité de matière de dihydrogène recueilli.
    \item Tout le zinc ayant disparu, déterminer la quantité de matière de zinc
          contenue dans l'échantillon.
    \item En déduire la composition massique (pourcentage en masse) du laiton.
    \item Calculer, pour la solution obtenue en fin de réaction, le
          concentration des ion \ce{H3O^{+}} et celle des ions
          \ch{Zn^{2+}_{(aq)}}.
\end{enumerate}

